\documentclass[aspectratio=169]{beamer}
\usetheme{Madrid}
\usecolortheme{whale}
\usepackage[utf8]{inputenc}
\usepackage[spanish]{babel}
\usepackage{amsmath,amsfonts,amssymb}
\usepackage{graphicx}
\usepackage{booktabs}
\usepackage{listings}
\usepackage{xcolor}

\lstset{
  basicstyle=\ttfamily\small,
  backgroundcolor=\color{gray!10},
  frame=single,
  breaklines=true,
  keywordstyle=\color{blue}\bfseries,
  commentstyle=\color{green!50!black},
  stringstyle=\color{red!70!black},
  showstringspaces=false
}

\title{Métodos de Análisis de Datos 1}
\subtitle{Clase 8: Estadística Descriptiva Univariada}
\author{Depto.\ de Matemática -- UNS}
\institute{Universidad Nacional del Sur}
\date{1er Cuatrimestre 2026}

\AtBeginSection[]{
  \begin{frame}{Contenidos}
    \tableofcontents[currentsection]
  \end{frame}
}

\begin{document}

\begin{frame}
  \titlepage
\end{frame}

\begin{frame}{Contenidos}
  \tableofcontents
\end{frame}

%=============================================================
\section{Introducción}
%=============================================================

\begin{frame}{¿Qué es la estadística descriptiva?}
  \begin{block}{Objetivo}
    Resumir y describir las características principales de un conjunto de datos, sin extrapolaciones ni inferencias sobre la población.
  \end{block}
  \bigskip
  \begin{itemize}
    \item Es el primer paso en cualquier análisis de datos.
    \item Permite detectar patrones, valores atípicos y distribuciones.
    \item Se aplica \textbf{variable por variable} (análisis univariado).
  \end{itemize}
  \bigskip
  \begin{alertblock}{Distinción clave}
    Estadística \textbf{descriptiva}: resume los datos observados. \quad
    Estadística \textbf{inferencial}: generaliza a la población.
  \end{alertblock}
\end{frame}

\begin{frame}{Tipos de variables}
  \begin{columns}
    \begin{column}{0.5\textwidth}
      \textbf{Variables cualitativas (categóricas)}
      \begin{itemize}
        \item \emph{Nominales}: sin orden. \\
          Ej.: sexo, grupo sanguíneo.
        \item \emph{Ordinales}: con orden. \\
          Ej.: nivel educativo, estadio clínico.
      \end{itemize}
    \end{column}
    \begin{column}{0.5\textwidth}
      \textbf{Variables cuantitativas (numéricas)}
      \begin{itemize}
        \item \emph{Discretas}: valores contables. \\
          Ej.: número de hijos.
        \item \emph{Continuas}: cualquier valor en un rango. \\
          Ej.: presión arterial, edad exacta.
      \end{itemize}
    \end{column}
  \end{columns}
  \bigskip
  \begin{block}{Importancia}
    El tipo de variable determina qué medidas y gráficos son adecuados.
  \end{block}
\end{frame}

%=============================================================
\section{Medidas de posición central}
%=============================================================

\begin{frame}{Media aritmética}
  Dado un conjunto de $n$ observaciones $x_1, x_2, \ldots, x_n$:
  \[
    \bar{x} = \frac{1}{n}\sum_{i=1}^{n} x_i
  \]
  \begin{itemize}
    \item Es la medida de tendencia central más usada.
    \item \textbf{Sensible} a valores extremos (outliers).
    \item Para datos agrupados: $\bar{x} = \frac{\sum f_i m_i}{\sum f_i}$, donde $m_i$ es la marca de clase.
  \end{itemize}
  \bigskip
  \begin{exampleblock}{Ejemplo}
    Edades: 25, 30, 28, 35, 22. \quad $\bar{x} = \frac{140}{5} = 28$ años.
  \end{exampleblock}
\end{frame}

\begin{frame}{Mediana y Moda}
  \begin{block}{Mediana ($Me$)}
    Valor que divide la distribución ordenada en dos mitades iguales.
    \begin{itemize}
      \item $n$ impar: $Me = x_{(n+1)/2}$
      \item $n$ par: $Me = \dfrac{x_{n/2} + x_{n/2+1}}{2}$
    \end{itemize}
    \textbf{Ventaja}: robusta frente a outliers.
  \end{block}
  \bigskip
  \begin{block}{Moda ($Mo$)}
    Valor (o clase) que aparece con mayor frecuencia.
    \begin{itemize}
      \item Puede no ser única: distribuciones bimodales o multimodales.
      \item Única medida aplicable a variables nominales.
    \end{itemize}
  \end{block}
\end{frame}

\begin{frame}{Cuantiles}
  Los \textbf{cuantiles} dividen la distribución ordenada en partes iguales:
  \bigskip
  \begin{itemize}
    \item \textbf{Cuartiles} ($Q_1, Q_2, Q_3$): dividen en 4 partes. $Q_2 = Me$.
    \item \textbf{Deciles} ($D_1,\ldots,D_9$): dividen en 10 partes.
    \item \textbf{Percentiles} ($P_1,\ldots,P_{99}$): dividen en 100 partes.
  \end{itemize}
  \bigskip
  \begin{alertblock}{Método de cálculo}
    Para el percentil $p$: calcular $L = p/100 \cdot n$. Si $L$ es entero, $P_p = (x_{(L)} + x_{(L+1)})/2$; si no, $P_p = x_{(\lceil L \rceil)}$.
  \end{alertblock}
\end{frame}

%=============================================================
\section{Medidas de dispersión}
%=============================================================

\begin{frame}{Varianza y desvío estándar}
  \begin{block}{Varianza muestral}
    \[
      s^2 = \frac{1}{n-1}\sum_{i=1}^{n}(x_i - \bar{x})^2
    \]
    Divisor $n-1$: corrección de Bessel (estimador insesgado).
  \end{block}
  \begin{block}{Desvío estándar (DE)}
    \[
      s = \sqrt{s^2}
    \]
    Tiene las \textbf{mismas unidades} que la variable original.
  \end{block}
  \begin{alertblock}{Sensibilidad}
    La varianza y el DE son muy sensibles a valores extremos.
  \end{alertblock}
\end{frame}

\begin{frame}{Rango, RIC y CV}
  \begin{block}{Rango (recorrido)}
    $R = x_{\max} - x_{\min}$ \quad — Muy sensible a outliers.
  \end{block}
  \bigskip
  \begin{block}{Rango Intercuartílico (RIC / IQR)}
    \[
      \text{RIC} = Q_3 - Q_1
    \]
    Contiene el 50\% central de los datos. \textbf{Robusto} frente a outliers.
  \end{block}
  \bigskip
  \begin{block}{Coeficiente de Variación (CV)}
    \[
      CV = \frac{s}{\bar{x}} \times 100\%
    \]
    Medida adimensional. Permite comparar dispersiones entre variables con distintas unidades o escalas.
  \end{block}
\end{frame}

\begin{frame}{Comparación de medidas de dispersión}
  \begin{center}
    \begin{tabular}{lccc}
      \toprule
      \textbf{Medida} & \textbf{Robusta} & \textbf{Unidades} & \textbf{Uso típico} \\
      \midrule
      Rango & No & Originales & Exploración rápida \\
      Varianza ($s^2$) & No & Cuadradas & Cálculos teóricos \\
      Desvío estándar ($s$) & No & Originales & Informe general \\
      RIC & Sí & Originales & Datos con outliers \\
      CV & No & Adimensional & Comparación cruzada \\
      \bottomrule
    \end{tabular}
  \end{center}
\end{frame}

%=============================================================
\section{Medidas de forma}
%=============================================================

\begin{frame}{Asimetría (Skewness)}
  \begin{block}{Coeficiente de asimetría de Fisher ($g_1$)}
    \[
      g_1 = \frac{1}{n}\frac{\sum_{i=1}^n (x_i - \bar{x})^3}{s^3}
    \]
  \end{block}
  \begin{itemize}
    \item $g_1 = 0$: distribución \textbf{simétrica}.
    \item $g_1 > 0$: cola a la \textbf{derecha} (asimetría positiva). $\bar{x} > Me > Mo$.
    \item $g_1 < 0$: cola a la \textbf{izquierda} (asimetría negativa). $\bar{x} < Me < Mo$.
  \end{itemize}
  \bigskip
  \begin{exampleblock}{Regla empírica}
    Si $|g_1| < 0.5$: aproximadamente simétrica. \quad $|g_1| > 1$: marcadamente asimétrica.
  \end{exampleblock}
\end{frame}

\begin{frame}{Curtosis (Kurtosis)}
  \begin{block}{Coeficiente de curtosis de Fisher ($g_2$)}
    \[
      g_2 = \frac{1}{n}\frac{\sum_{i=1}^n (x_i - \bar{x})^4}{s^4} - 3
    \]
    (exceso de curtosis respecto a la distribución Normal)
  \end{block}
  \begin{itemize}
    \item $g_2 = 0$: \textbf{Mesocúrtica} (similar a la Normal).
    \item $g_2 > 0$: \textbf{Leptocúrtica} — pico pronunciado, colas pesadas.
    \item $g_2 < 0$: \textbf{Platicúrtica} — pico aplanado, colas livianas.
  \end{itemize}
\end{frame}

%=============================================================
\section{Tablas de frecuencia}
%=============================================================

\begin{frame}{Tabla de frecuencias para datos cualitativos}
  \begin{center}
    \begin{tabular}{lccc}
      \toprule
      \textbf{Categoría} & \textbf{Frec.\ abs.} ($n_i$) & \textbf{Frec.\ rel.} ($f_i$) & \textbf{Frec.\ acum.} ($F_i$) \\
      \midrule
      A & 15 & 0.30 & 0.30 \\
      B & 22 & 0.44 & 0.74 \\
      C & 13 & 0.26 & 1.00 \\
      \midrule
      \textbf{Total} & 50 & 1.00 & --- \\
      \bottomrule
    \end{tabular}
  \end{center}
  \bigskip
  \begin{itemize}
    \item $f_i = n_i / n$: proporción de cada categoría.
    \item $F_i = \sum_{j \le i} f_j$: frecuencia acumulada (para ordinales o cuantitativas).
  \end{itemize}
\end{frame}

\begin{frame}{Tabla de frecuencias para datos cuantitativos}
  Para variables continuas se agrupan en \textbf{clases} (intervalos):
  \begin{itemize}
    \item Número de clases recomendado: regla de Sturges $k = 1 + 3.322 \log_{10}(n)$.
    \item Amplitud de clase: $h = (x_{\max} - x_{\min})/k$.
    \item Marca de clase: $m_i = (L_i + L_{i+1})/2$.
  \end{itemize}
  \bigskip
  \begin{alertblock}{Convención}
    Los intervalos son cerrados a izquierda y abiertos a derecha: $[L_i, L_{i+1})$, excepto el último que se cierra en ambos extremos.
  \end{alertblock}
\end{frame}

%=============================================================
\section{Resumen}
%=============================================================

\begin{frame}{Resumen de la clase}
  \begin{block}{Medidas estudiadas}
    \begin{itemize}
      \item \textbf{Posición}: media, mediana, moda, cuantiles.
      \item \textbf{Dispersión}: rango, varianza, DE, RIC, CV.
      \item \textbf{Forma}: asimetría ($g_1$), curtosis ($g_2$).
    \end{itemize}
  \end{block}
  \bigskip
  \begin{block}{Próxima clase}
    Clase 9: Análisis de distribuciones — histogramas, boxplots, Q-Q plots y ajuste de distribuciones teóricas.
  \end{block}
\end{frame}

\end{document}
